\section{A global localized theorem for the \texorpdfstring{\(\rh/\cB_1\)}{rho/B1} component}
\label{sec:global-rho}

The three finite-prime restorations in \cref{sec:137} show that the
rank losses inherited from the control-prime presentation are not
intrinsic at those primes.  They do not, by themselves, exclude all
other primes.  We now replace that finite screen by an exact
characteristic-zero certificate for the same complete component map.

Put
\[
  R=\ZZ[\zeta],\qquad \zeta=\zeta_{17},\qquad
  R_S=R\left[\frac1{2\cdot3\cdot17}\right].
\]
The power basis
\[
  1,\zeta,\ldots,\zeta^{15}
\]
is an integral basis of the cyclotomic ring \(R\)
\cite{Washington1997}.  In the fixed bases used throughout the component
audit, the complete minimum-circuit map is
\[
  A=F_{\rh,1}^{\mathrm{full}}:R^{1317}\longrightarrow R^{1312}.
\]
Thus its cokernel measures only the residual
\(\rh/\cB_1\) component obstruction.  It is not the complete group
\(\Theta_1\).

\subsection{Exact lift and determinant reconstruction}

\begin{lemma}[Common cyclotomic lift]
\label{lem:common-cyclotomic-lift}
The \(1\,317\) source columns and every complete minimum-circuit equation
for \(A\) have one exact realization over \(R\) in the power basis above.
Under a residue homomorphism sending \(\zeta\) to an element of order
\(17\), this realization reduces entrywise to the corresponding complete
finite component map.
\end{lemma}

\begin{proof}
The fixed-space columns are first normalized as exponent vectors modulo
\(17\).  Every representation entry is then a signed sum of those powers,
reduced by
\[
  1+\zeta+\cdots+\zeta^{16}=0.
\]
Circuit rows are formed from the same six signed block columns before any
specialization.  Consequently evaluation at a seventeenth root commutes
with both column construction and row formation.  The independent replay
checks all \(1\,317\) normalized columns and the entrywise reductions used
at the three screened primes; no independently chosen finite-field basis
is lifted back to characteristic zero.
\end{proof}

For completeness, we record why the determinant lifts below are exact
rather than heuristic modular reconstructions.

\begin{lemma}[Certified signed CRT reconstruction]
\label{lem:certified-crt}
Let \(B\) be either square matrix in \cref{prop:two-exact-minors}, and
write
\[
  \det B=\sum_{j=0}^{15}c_j\zeta^j.
\]
For each complex embedding, Hadamard's inequality bounds the determinant
by the square root of the product of the certified row energies.  Fourier
inversion in the seventeenth cyclotomic basis then gives
\[
 \max_j |c_j|
 \ \leq\ \frac{32}{17}
 \max_{\sigma:R\hookrightarrow\CC}|\sigma(\det B)|.
\]
If the product \(M\) of the split CRT primes is larger than twice this
coefficient bound, the sixteen signed coefficients \(c_j\) are uniquely
determined by their residues modulo \(M\).
\end{lemma}

\begin{proof}
Hadamard gives
\(
 |\det\sigma(B)|^2\leq\prod_r E_r
\)
from the row-energy bounds \(E_r\).  Applying the inverse
seventeenth-root transform and eliminating the coefficient of
\(\zeta^{16}\) with \(\Phi_{17}(\zeta)=0\) gives the displayed
\(32/17\) bound.  Finally, two integers in the signed interval
\((-M/2,M/2)\) with the same residue modulo \(M\) are equal.  Apply this
coordinatewise.
\end{proof}

\begin{proposition}[Two exact maximal minors]
\label{prop:two-exact-minors}
The matrix \(A\) has two explicitly displayed
\(1312\times1312\) minors with determinants
\(\delta_{\mathrm G},\delta_{\mathrm R}\in R\).  Their independent exact
reconstructions have the following certified data.
\[
\begin{array}{lrrrrrr}
\toprule
\text{minor}&\text{entries}&\text{check}&\text{CRT primes}&
\text{bound bits}&\text{modulus bits}&\text{norm bits}\\
\midrule
\delta_{\mathrm G}&7830&89\bmod103&96&2937&2967&4833\\
\delta_{\mathrm R}&7819&107\bmod137&93&2872&2874&2003\\
\bottomrule
\end{array}
\]
The largest reconstructed coefficient has respectively \(304\) and
\(128\) bits.
\end{proposition}

\begin{proof}
For \(\delta_{\mathrm G}\), the complete map at
\((103,\zeta-8)\) has rank \(1312\).  Replacing the unique deficient
column in the original selected minor gives determinant \(89\) and an
exact sparse lift with \(7\,830\) nonzero entries.

For \(\delta_{\mathrm R}\), scan the \(11\,274\) circuit orbits in reverse
canonical order and, inside an orbit, scan the four representation
coordinates in reverse order.  The first full row basis is reached after
\(11\,069\) orbits and \(44\,276\) rows.  With canonical columns
\(0,\ldots,1311\), its determinant is \(107\) at
\((137,\zeta-122)\), and its exact lift has \(7\,819\) nonzero entries.

In both cases an independent binary parser reconstructs the transpose,
checks the finite determinant, and obtains rank \(1311\) after the
designated deletion or duplication mutation.  The split-prime
congruences, row-energy products, and inequalities in
\cref{lem:certified-crt} are then replayed for every CRT prime.  Since the
recorded moduli exceed twice the proved bounds, the lifted coefficient
vectors, and hence their exact cyclotomic norms, are unique.
\end{proof}

\subsection{The two-minor ideal certificate}

We use two standard consequences of Fitting ideals and integral Smith
theory \cite{StacksProjectFitting,Cohen1993,Newman1972}.

\begin{lemma}[Two-minor Fitting criterion]
\label{lem:two-minor-fitting}
Let \(T\) be a commutative ring and
\(f:T^n\to T^m\), with \(n\geq m\).  If two maximal minors
\(\alpha,\beta\) of \(f\) satisfy
\((\alpha,\beta)T'=T'\) after a flat base change \(T\to T'\), then
\[
  \coker(f)\otimes_TT'=0.
\]
\end{lemma}

\begin{proof}
The zeroth Fitting ideal of \(\coker(f)\) is generated by all maximal
minors of a presentation matrix and commutes with base change.  It
therefore contains \((\alpha,\beta)T'=T'\).  Thus its localization at
every prime of \(T'\) is the unit ideal.  Some maximal minor is a unit in
each local ring, so the localized map is surjective.  The finitely
generated cokernel has empty support and is zero.
\end{proof}

\begin{lemma}[Multiplication-lattice certificate]
\label{lem:multiplication-lattice}
For \(\alpha\in R\), let \(M(\alpha)\in\operatorname{Mat}_{16}(\ZZ)\)
be multiplication by \(\alpha\) in the integral power basis.  Then
\[
 \im_{\ZZ}M(\alpha)=\alpha R.
\]
Consequently the columns of
\[
 \bigl[M(\alpha)\ M(\beta)\bigr]\in
 \operatorname{Mat}_{16\times32}(\ZZ)
\]
span \(\alpha R+\beta R\), and the nonzero Smith entries of this
integer matrix are the additive invariant factors of
\(R/(\alpha R+\beta R)\).
\end{lemma}

\begin{proof}
The columns of \(M(\alpha)\) are
\(\alpha,\alpha\zeta,\ldots,\alpha\zeta^{15}\) in coordinates, hence
span exactly the principal ideal \(\alpha R\).  Concatenating two such
matrices adds their column lattices.  Smith normal form computes the
quotient of the ambient free abelian group by that full-rank sublattice.
\end{proof}

\begin{theorem}[Global localized saturation of the
\(\rh/\cB_1\) component]
\label{thm:global-rho}
For the complete minimum-circuit equation map
\[
 A=F_{\rh,1}^{\mathrm{full}}:R^{1317}\longrightarrow R^{1312},
\]
one has
\[
 \coker(A)\otimes_RR_S=0.
\]
Equivalently, \(A\otimes_RR_S\) is surjective.  Thus the complete
\(\rh/\cB_1\) component cokernel has no prime-ideal support away from the
rational primes \(2\), \(3\), and \(17\).
\end{theorem}

\begin{proof}
Apply \cref{lem:multiplication-lattice} to the two determinants in
\cref{prop:two-exact-minors}.  The independently replayed Smith form of
\[
 \bigl[M(\delta_{\mathrm G})\ M(\delta_{\mathrm R})\bigr]
\]
has diagonal
\[
 \underbrace{17,\ldots,17}_{11\ \mathrm{times}},
 \qquad
 \underbrace{289,\ldots,289}_{5\ \mathrm{times}}.
\]
Therefore
\[
 [R:\delta_{\mathrm G}R+\delta_{\mathrm R}R]
 =17^{11}(17^2)^5
 =17^{21}
 =69\,091\,933\,913\,008\,732\,880\,827\,217.
\]
After inverting \(17\), and hence over \(R_S\), the two determinants
generate the unit ideal.  Both are maximal minors of \(A\), so
\cref{lem:two-minor-fitting} proves the assertion.
\end{proof}

\begin{remark}[Why a stable determinant factor is not support]
\label{rem:minor-artifact}
Neither displayed determinant is an \(S\)-unit by itself.  Moreover,
several one-column replacements and a reverse-column basis built on the
same frozen row set retain a large common non-\(S\) factor.  That
persistence is a presentation artifact: the genuinely different
reverse-row minor removes it from the ideal sum.

As a negative control, replacing \(\delta_{\mathrm G}\) by the original
G0 determinant leaves, after removal of \(2\), \(3\), and \(17\), the
nonunit residual
\[
 47\,821\,880\,003\,927\,349\,029.
\]
Thus neither a factor of one norm nor its survival across a closely
related minor family can be promoted to support of the complete
cokernel.
\end{remark}

\begin{remark}[Exact boundary]
\label{rem:global-rho-boundary}
\Cref{thm:global-rho} closes one of the six large component maps.  The
global gates for
\[
 \rh\om/\cB_1,\quad
 \rh/\cB_2,\quad
 \rh\om/\cB_2,\quad
 \rh/\cB_3,\quad
 \rh\om/\cB_3
\]
remain open.  In particular, the theorem does not prove
\(\operatorname{Supp}(\Theta_i)\subseteq\{2,3,17\}\), does not determine
the global Smith groups, and is not a tomography or Markov-connectivity
statement.
\end{remark}
